\begin{workedexample}[Worked Example: Maximum unilateral gain]
A device has $|S_{21}| = 4$, $|S_{11}| = 0.5$, $|S_{22}| = 0.4$ (and
$S_{12}\!\approx\!0$). Its maximum unilateral transducer gain is
\[ G_{TU,\max} = \frac{|S_{21}|^2}{(1-|S_{11}|^2)(1-|S_{22}|^2)}
   = \frac{16}{(1-0.25)(1-0.16)} = \frac{16}{0.63} = 25.4, \]
i.e.\ $14.0\ \text{dB}$, obtained by conjugately matching both ports.
\end{workedexample}

\begin{keytakeaways}
\begin{itemize}[leftmargin=1.4em,itemsep=2pt]
\item $G_{TU,\max}=|S_{21}|^2/[(1-|S_{11}|^2)(1-|S_{22}|^2)]$ with input/output conjugate match.
\item Check stability ($K>1$, $|\Delta|<1$) before choosing terminations.
\item Constant-gain circles trade gain for bandwidth, match, or noise.
\end{itemize}
\end{keytakeaways}

\begin{exercises}
\item $|S_{21}| = 3$, $|S_{11}| = 0.4$, $|S_{22}| = 0.3$. Find $G_{TU,\max}$ in dB.
\item Why can maximum available gain be undefined for a real device?
\item What does a constant-gain circle represent?
\end{exercises}

\furtherreading{Gonzalez~\cite{gonzalez} (ch.~3); Pozar~\cite{pozar} (ch.~12).}
